%%%%%%%%%%  scan idx 226  |  folio 213  %%%%%%%%%%
(3.0.2))~;
\item $\varphi=\varprojlim\varphi_r$ (i.e.
pour $u\in\Diff^n_\mathbb C(E,F)$, $\varphi(u)$ est la limite du système projectif des $\varphi_r(u)$).
\end{enumerate}
(Les assertions précédentes sont faciles à partir des résultats de (loc.cit)~: on se ramène aussitôt à $n=0$, (ii) découle alors directement des définitions, tandis que (i) résulte (après localisation du problème grâce à une partition de l'unité) de la formule qui donne les dérivées successives d'un produit).

\textbf{Lemme 3.1.4.} Soient $S$ un espace topologique, $V$ une variété compacte, $E$ et $F$ des fibrés vectoriels complexes de rang fini sur $V$, et $n\in\N$. Notons $X$ la $S$-variété mixte $S\times V$. Pour $u\in\Diff^n_\mathbb C(\pr_2^*(E),\pr_2^*(F))$, l'application $s\mapsto i_s^*(u)$ de $S$ dans $\Diff^n_\mathbb C(E,F)$ est continue. L'application linéaire
\[
  \Diff^n_\mathbb C(\pr_2^*(E),\pr_2^*(F))
  \longrightarrow C(S,\Diff^n_\mathbb C(E,F))
\]
ainsi définie est un isomorphisme.

\emph{Preuve.} Moyennant l'identification (1.2.3.1), on peut se borner à $n=0$. Posant $\underline{\Homop}(E,F)=G$, on est alors ramené à montrer que pour $u\in\Gamma(X,\pr_2^*(G))$, l'application $s\mapsto u_s=i_s^*(u)$ de $S$ dans $C^\infty(V,G)$ est continue, et que l'application linéaire
\[
  \Gamma(X,\pr_2^*(G))\longrightarrow C(S,C^\infty(V,G))
\]
ainsi obtenue est un isomorphisme. Mais cela résulte immédiatement de la définition des variétés mixtes (1.1).

\textbf{3.1.5.} Sous les hypothèses de (3.1.4), soit
\[
  u\in\Diff^n(\pr_2^*(E),\pr_2^*(F)).\text{\footnotemark}
\]
\footnotetext{On omet l'indice $\mathbb C$ pour abréger.}
